% GENERATED FILE. Edit canonical lessons, then run npm run generate:books.
% canonical-source-hash: fnv1a64:cda414d9f0479498
% canonical-lessons: UR-C294-qahat, UR-C294-faqa, UR-C294-vaba, UR-C294-sarna, UR-C294-haiza

\chapter{Famine and Sickness}
\label{ch:ur-famine-and-sickness}
\hlchaptermodality{294}

\begin{chapteropening}
\textbf{By the end of this chapter:} \emph{I can tell how hunger and disease spread.}

Complete the last lesson of chapter 294: I can tell how hunger and disease spread.
\end{chapteropening}

\section[\texorpdfstring{\ur{قحط} (qahat)}{قحط (qahat)}]{\ur{قحط} (qahat) — a famine}

\label{lesson:UR-C294-qahat}



\begin{quote}
Before the new one: say the Urdu for a wave, then the Urdu for a whirlpool.
\end{quote}

\subsection*{You'll want to know: \ur{قحط}}
\textbf{\ur{قحط}} — \emph{qahat} — \textquotedblleft{}a famine\textquotedblright{}.

\begin{quote}
\ur{اس} \ur{سال} \ur{گاؤں} \ur{میں} \ur{قحط} \ur{پڑا۔}
\end{quote}

\begin{quote}
\emph{is sāl gāoṅ meṅ qahat paṛā.}
\end{quote}

\begin{quote}
This year there was a famine in the village.
\end{quote}

Masculine.

\begin{cousinweb}
\textbf{\ur{قحط}} (\emph{qahat}) is usually traced, through Persian, to Arabic \emph{qaḥṭ}, \textquotedblleft{}drought, a failure of rain\textquotedblright{}. A famine is said to fall: \textbf{\ur{قحط} \ur{پڑا}}.
\end{cousinweb}

\begin{grammarlens}[title={What you've built}]
One more word for telling what happened.
\end{grammarlens}

\subsection*{Your turn}
\begin{itemize}
\raggedright
  \item \emph{Say it:} \emph{qahat}
  \item \emph{Say it:} \emph{qahat}, once more
  \item \emph{Say it:} \emph{is sāl gāoṅ meṅ qahat paṛā.} — This year there was a famine in the village.
  \item \emph{Recall:} say the Urdu for a wave, then the Urdu for a whirlpool, then say \emph{qahat} again
\end{itemize}

\begin{tcolorbox}[breakable,colback=teal!4,colframe=teal!35!black,title={Before you move on}]
What does \ur{قحط} mean? (\textquotedblleft{}A famine\textquotedblright{}.) Say it once more.
\end{tcolorbox}
\section[\texorpdfstring{\ur{فاقہ} (fāqā)}{فاقہ (fāqā)}]{\ur{فاقہ} (fāqā) — going without food}

\label{lesson:UR-C294-faqa}



\begin{quote}
Before the new one: say the Urdu for a whirlpool, then the Urdu for a famine.
\end{quote}

\subsection*{You'll want to know: \ur{فاقہ}}
\textbf{\ur{فاقہ}} — \emph{fāqā} — \textquotedblleft{}going without food\textquotedblright{}.

\begin{quote}
\ur{لوگ} \ur{فاقے} \ur{سے} \ur{کمزور} \ur{ہو} \ur{گئے۔}
\end{quote}

\begin{quote}
\emph{log fāqe se kamzor ho gae.}
\end{quote}

\begin{quote}
People grew weak from hunger.
\end{quote}

Masculine.

\begin{cousinweb}
is usually traced, through Persian, to Arabic \emph{fāqa}, \textquotedblleft{}want, need\textquotedblright{}. A \textbf{\ur{فاقہ}} is a day with nothing to eat.
\end{cousinweb}

\begin{grammarlens}[title={What you've built}]
One more word for telling what happened.
\end{grammarlens}

\subsection*{Your turn}
\begin{itemize}
\raggedright
  \item \emph{Say it:} \emph{fāqā}
  \item \emph{Say it:} \emph{fāqā}, once more
  \item \emph{Say it:} \emph{log fāqe se kamzor ho gae.} — People grew weak from hunger.
  \item \emph{Recall:} say the Urdu for a whirlpool, then the Urdu for a famine, then say \emph{fāqā} again
\end{itemize}

\begin{tcolorbox}[breakable,colback=teal!4,colframe=teal!35!black,title={Before you move on}]
What does \ur{فاقہ} mean? (\textquotedblleft{}Going without food\textquotedblright{}.) Say it once more.
\end{tcolorbox}
\section[\texorpdfstring{\ur{وبا} (vabā)}{وبا (vabā)}]{\ur{وبا} (vabā) — an epidemic}

\label{lesson:UR-C294-vaba}



\begin{quote}
Before the new one: say the Urdu for a famine, then the Urdu for going without food.
\end{quote}

\subsection*{You'll want to know: \ur{وبا}}
\textbf{\ur{وبا}} — \emph{vabā} — \textquotedblleft{}an epidemic\textquotedblright{}.

\begin{quote}
\ur{لوگوں} \ur{کو} \ur{وبا} \ur{کا} \ur{ڈر} \ur{تھا۔}
\end{quote}

\begin{quote}
\emph{logoṅ ko vabā kā ḍar thā.}
\end{quote}

\begin{quote}
People were afraid of the epidemic.
\end{quote}

Feminine.

\begin{cousinweb}
\textbf{\ur{وبا}} (\emph{vabā}) is usually traced, through Persian, to Arabic \emph{wabā}, \textquotedblleft{}a plague\textquotedblright{}. It is feminine.
\end{cousinweb}

\begin{grammarlens}[title={What you've built}]
One more word for telling what happened.
\end{grammarlens}

\subsection*{Your turn}
\begin{itemize}
\raggedright
  \item \emph{Say it:} \emph{vabā}
  \item \emph{Say it:} \emph{vabā}, once more
  \item \emph{Say it:} \emph{logoṅ ko vabā kā ḍar thā.} — People were afraid of the epidemic.
  \item \emph{Recall:} say the Urdu for a famine, then the Urdu for going without food, then say \emph{vabā} again
\end{itemize}

\begin{tcolorbox}[breakable,colback=teal!4,colframe=teal!35!black,title={Before you move on}]
What does \ur{وبا} mean? (\textquotedblleft{}An epidemic\textquotedblright{}.) Say it once more.
\end{tcolorbox}
\section[\texorpdfstring{\ur{سڑنا} (saṛnā)}{سڑنا (saṛnā)}]{\ur{سڑنا} (saṛnā) — to rot}

\label{lesson:UR-C294-sarna}



\begin{quote}
Before the new one: say the Urdu for going without food, then the Urdu for an epidemic.
\end{quote}

\subsection*{You'll want to know: \ur{سڑنا}}
\textbf{\ur{سڑنا}} — \emph{saṛnā} — \textquotedblleft{}to rot\textquotedblright{}.

\begin{quote}
\ur{گرمی} \ur{سے} \ur{پھل} \ur{سڑ} \ur{گئے۔}
\end{quote}

\begin{quote}
\emph{garmī se phal saṛ gae.}
\end{quote}

\begin{quote}
The fruit rotted in the heat.
\end{quote}

Past (it …): \ur{وہ} \ur{پھل} \ur{سڑ} \ur{گیا}, \emph{voh phal saṛ gayā}.

\begin{cousinweb}
\textbf{\ur{سڑنا}} (\emph{saṛnā}) is inherited: CDIAL lists the Hindi cousin \emph{saṛnā}, \textquotedblleft{}to rot\textquotedblright{}, under Sanskrit \emph{śáṭati} (CDIAL 12268).
\end{cousinweb}

\begin{grammarlens}[title={What you've built}]
One more word for telling what happened.
\end{grammarlens}

\subsection*{Your turn}
\begin{itemize}
\raggedright
  \item \emph{Say it:} \emph{saṛnā}
  \item \emph{Say it:} \emph{saṛnā}, once more
  \item \emph{Say it:} \emph{garmī se phal saṛ gae.} — The fruit rotted in the heat.
  \item \emph{Recall:} say the Urdu for going without food, then the Urdu for an epidemic, then say \emph{saṛnā} again
\end{itemize}

\begin{tcolorbox}[breakable,colback=teal!4,colframe=teal!35!black,title={Before you move on}]
What does \ur{سڑنا} mean? (\textquotedblleft{}To rot\textquotedblright{}.) Say it once more.
\end{tcolorbox}
\section[\texorpdfstring{\ur{ہیضہ} (haizā)}{ہیضہ (haizā)}]{\ur{ہیضہ} (haizā) — cholera}

\label{lesson:UR-C294-haiza}



\begin{quote}
Before the new one: say the Urdu for an epidemic, then the Urdu for to rot.
\end{quote}

\subsection*{You'll want to know: \ur{ہیضہ}}
\textbf{\ur{ہیضہ}} — \emph{haizā} — \textquotedblleft{}cholera\textquotedblright{}.

\begin{quote}
\ur{ہیضہ} \ur{ایک} \ur{خطرناک} \ur{بیماری} \ur{ہے۔}
\end{quote}

\begin{quote}
\emph{haizā ek khatarnāk bīmārī hai.}
\end{quote}

\begin{quote}
Cholera is a dangerous illness.
\end{quote}

Masculine.

\begin{cousinweb}
\textbf{\ur{ہیضہ}} (\emph{haizā}) is usually traced, through Persian, to Arabic \emph{hayḍa}, \textquotedblleft{}cholera\textquotedblright{}. It spread where water was dirty.
\end{cousinweb}

\begin{grammarlens}[title={What you've built}]
One more word for telling what happened.
\end{grammarlens}

\subsection*{Your turn}
\begin{itemize}
\raggedright
  \item \emph{Say it:} \emph{haizā}
  \item \emph{Say it:} \emph{haizā}, once more
  \item \emph{Say it:} \emph{haizā ek khatarnāk bīmārī hai.} — Cholera is a dangerous illness.
  \item \emph{Recall:} say the Urdu for an epidemic, then the Urdu for to rot, then say \emph{haizā} again
\end{itemize}

\begin{tcolorbox}[breakable,colback=teal!4,colframe=teal!35!black,title={Before you move on}]
What does \ur{ہیضہ} mean? (\textquotedblleft{}Cholera\textquotedblright{}.) Say it once more.
\end{tcolorbox}
